\section{Complex low fusion products and their fourth moments}
\label{sec:complex-low-fourth}

The remaining non-self-dual quadratic tensor $\RR{4}{155}$ illustrates why
survival certificates should be retained as carefully as exclusion
certificates. Its invertible group is $C_2\times C_2$, its two noninvertible
simples are dual to one another, and their common Frobenius--Perron dimension
is $d=4+3\sqrt2$. Thus
\[
 \operatorname{FPdim}(\mathcal C)=4+2d^2=72+48\sqrt2,
 \qquad p_D(t)=t^2-144t+576.
\]
The criteria below use the entire input tensor and are not specific to this
example. All categorical conclusions are conditional on a pseudo-unitary
categorification; the accompanying signed-reduction certificate separately
checks the passage from unrestricted categorification to a pseudo-unitary
Galois conjugate for this tensor.

\subsection{Complex low Verlinde coefficients}
Write $D=\operatorname{FPdim}(\mathcal C)$, and let $A_a$ denote the simple
summands of $I(\mathbf1)$, with $A_0=\mathbf1$. For a central simple $Z$ set
$d_Z=\operatorname{FPdim}(Z)$ and
\[
 \lambda_a(Z)=\frac{S_{A_a,Z}}{S_{\mathbf1,Z}}.
\]
The induction character table and the restriction column of $Z$ determine
these numbers. They need not be real. Complex conjugation of the complete
character table determines the permutation $a\mapsto a^*$, and one must
retain this permutation when writing restriction-product equations.
The Verlinde formula gives
\begin{equation}\label{eq:r4-complex-triple}
 t_{abc}:=\dim\operatorname{Hom}(\mathbf1,A_aA_bA_c)
 =\frac1{D^2}\sum_Z
       d_Z^2\lambda_a(Z)\lambda_b(Z)\lambda_c(Z)
 \in\mathbb Z_{\ge0}.
\end{equation}
In particular $N_{A_a,A_b}^{A_c}=t_{ab c^*}$, not necessarily $t_{abc}$.
This identity is linear in the numbers of high central simples with each
restriction column and relative twist-trace type.

For a high restriction column $v$, put
\[
 x_v=\#\{Z:F(Z)=v\},\qquad
 q_{ab,v}=\sum_{F(Z)=v}N_{A_a,A_b}^Z.
\]
If $f_{ab}$ is the coefficient vector of $F(A_a)F(A_b)$, tensoriality of $F$
gives the exact componentwise equality
\begin{equation}\label{eq:r4-complex-restriction}
 f_{ab}=\delta_{a,b^*}e_0+
          \sum_{c\ne0}t_{ab c^*}F(A_c)+\sum_v q_{ab,v}v.
\end{equation}
For every nonzero coordinate of $v$, nonnegativity gives
$N_{A_a,A_b}^Z v_i\le(f_{ab})_i$. Consequently, with
\[
 B_{ab,v}=\min_{i:v_i>0}\left\lfloor\frac{(f_{ab})_i}{v_i}\right\rfloor,
 \qquad 0\le q_{ab,v}\le B_{ab,v}x_v.
\]
All variables in these equations are nonnegative integers in an actual
category. Linear infeasibility is therefore an obstruction; a solution,
even integral, constructs no category.

\subsection{Fourth moments and integer squares}
There is a useful further identity that does not require the input tensor
to be self-dual:
\begin{equation}\label{eq:r4-fourth-moment}
 Q_{ab}:=\frac1{D^2}\sum_Z d_Z^2
  \lambda_a(Z)\lambda_b(Z)\lambda_{a^*}(Z)\lambda_{b^*}(Z)
 =\sum_Z\bigl(N_{A_a,A_b}^Z\bigr)^2.
\end{equation}
Indeed the left side is the coefficient of the unit in
$(A_aA_b)(A_aA_b)^*$, by simultaneous diagonalization of the fusion matrices.
Expanding $A_aA_b$ in the simple basis gives the right side by rigidity.
Thus this assertion uses the based-ring involution and positivity of fusion
multiplicities, not an assumption that the category is unitary.

Introduce $z_{ab,v}=\sum_{F(Z)=v}(N_{A_a,A_b}^Z)^2$ and corresponding
square variables for the known low coefficients. For integers $n\ge0$ and
$0\le n\le B$, the identities
\[
 n^2-(2k+1)n+k(k+1)=(n-k)(n-k-1)\ge0,
 \qquad n^2\le Bn
\]
hold for every integer $k\ge0$. Summing them over a restriction type gives
\begin{equation}\label{eq:r4-square-cuts}
 (2k+1)q_{ab,v}-k(k+1)x_v\le z_{ab,v}
          \le B_{ab,v}q_{ab,v}.
\end{equation}
For a low simple, replace $x_v$ by $1$ and $q_{ab,v}$ by its individual
Verlinde multiplicity. Equation~\eqref{eq:r4-fourth-moment} becomes the sum
of these square variables, together with the unit contribution
$\delta_{a,b^*}$. Any finite selection of the lower inequalities is valid;
the implementation supports both all $0\le k<B$ and a stated cutoff.
This gives an exact linear strengthening, with integrality retained in
bounded discovery experiments. Failure of a mixed-integer optimizer to
finish is recorded only as an incomplete search.

\subsection{Transport by central invertibles}
Suppose the low frame has exactly one column equal to $e_0$.
A central invertible restricting to the unit must occur in $I(\mathbf1)$,
so the forgetful homomorphism on invertible objects has trivial kernel.
Its image is therefore a subgroup $K$ of the based-ring invertible group,
and each $g\in K$ has exactly one invertible lift.
For any such lift $L$, tensoring by $L$ permutes all central simples and sends
a restriction vector $v$ to the coefficient vector of $gv$. In particular,
the total counts of these restriction types agree, including low objects.

There is a stronger constraint when $g^2=1$ and $\nu_2(g)=-1$.
The lift has order two and twist $\theta_L=\pm i$. Its monodromy with a
simple $Z$ is a scalar in $\{1,-1\}$: monodromy with invertibles is
multiplicative, and $L^2=\mathbf1$. The balancing identity therefore gives
\[
 \theta_{LZ}=\pm i\theta_Z.
\]
In particular $LZ\not\cong Z$, so this permutation is a fixed-point-free
involution. When the ambient cyclotomic field contains $i$, normalized
relative traces transform by
$(t_j)_j\mapsto((\pm i)^j t_j)_j$.
One can encode the resulting matching with nonnegative integral edge
variables between allowed pairs of restriction/trace types. An edge
joining distinct types contributes one to each marginal; a loop in the
\emph{type} graph contributes two, because it still represents two distinct
simples. This last convention is necessary when relative trace loses the
underlying twist information. These matching equations are necessary only;
for several lifts they do not impose all simultaneous group-action
relations.

\subsection{Recorded scope and exact surviving branches}
For $\RR{4}{155}$ the signed-reduction replay reconstructs all $40$ signed
quotients in three sign orbits. Exactly one orbit consists of sign gauges
of the original tensor, and eight nongauge character components are
rejected. Thus every categorification is spherical and has a pseudo-unitary
Galois conjugate, by the signed-reduction theorem used earlier.

The complex low-Verlinde and restriction-product system admits an exact
nonnegative integral solution at ambient conductor $8$ and trace height
$6$: it has $2,674$ equations, $4,295$ variables, and total central-object
count $48$. Enlarging the ambient field to conductor $24$, at height $3$,
also admits an exact integral solution, with $2,940$ equations and $4,698$
variables. These solutions are replayed using exact rational arithmetic
with numerical optimizers disabled. They certify survival of precisely
the stated necessary systems. They are not modular data, and their
central-object counts are not asserted to be the rank of an existing
center.

The combined fourth-moment and central-invertible system also has an
exact integral survivor at conductor $8$, height $3$, and square-cut
cutoff $8$: $7,764$ equations in $9,982$ variables, with central-object
count $36$ and pointed image $\{1,h\}$. Its optimized-Python replay
checks every equation with numerical optimizers disabled. Their complete
per-input applicability records, including all ten input frontier rings
and all $54$ positive controls, are in
\texttt{research/round4/m4\_155/}. No timeout or numerical feasibility
result is used as a categorical decision. The exact witnesses are stored
separately from exclusion certificates.
